\chapter{Doelstelling en Onderzoeksvragen}

\section{Doelstelling}
% TIP:
% Formuleer dit iets praktischer dan bij afstuderen.
% Voorbeeld:
% "Het doel van dit stageonderzoek is om te onderzoeken en te ontwerpen hoe
% een koppeling tussen SolvCRM en e-Boekhouden.nl gerealiseerd kan worden,
% zodat boekhoudgegevens automatisch en foutloos kunnen worden uitgewisseld
% en Solvware-klanten tijd besparen."

Het doel van het dit stageonderzoek is om te onderzoeken en te ontwerpen hoe een koppeling tussen SolvCRM en e-boekhouden.nl gerealiseerd kan worden, boekhoudgegevens automatisch en foutloos kunnen worden uitgewisseld tussen de twee systemen, en Solvware klanten tijd besparen.

\section{Hoofdvraag}
% TIP: Neem de hoofdvraag exact op.
% "Hoe kan een koppeling tussen SolvCRM en e-Boekhouden.nl worden gerealiseerd
% zodat boekhoudgegevens automatisch en foutloos worden uitgewisseld,
% waardoor Solvware-klanten tijd besparen?"

De hoofdvraag van dit onderzoek luid:
"Hoe kan een koppeling tussen SolvCRM en e-Boekhouden.nl worden gerealiseerd
zodat boekhoudgegevens automatisch en foutloos worden uitgewisseld,
waardoor Solvware-klanten tijd besparen?"

\section{Deelvragen}
% TIP: Gewoon je deelvragen opnemen.
\begin{enumerate}
    \item Deelvraag 1: Hoe worden boekhoudgegevens tussen SolvCRM en e-Boekhouden.nl momenteel uitgewisseld?
    \item Deelvraag 2: Welke functionele en niet-functionele eisen zijn noodzakelijk voor het ontwerp van een koppeling tussen SolvCRM en e-Boekhouden.nl?
    \item Deelvraag 3: Welke technische oplossingsrichtingen bestaan er voor het automatiseren van de uitwisseling van klant- en boekhoudgegevens, en welke hiervan is het meest geschikt voor Solvware?
    \item Deelvraag 4: In welke mate vermindert de ontworpen koppeling de tijdsduur en foutgevoeligheid van de uitwisseling van boekhoudgegevens?
    \item Deelvraag 5: Welke best practices zijn nodig om een kwalitatieve Proof of Concept te bouwen, testen en onderhouden binnen Solvware?
\end{enumerate}
